\documentclass[aps,preprint]{revtex4}%
\usepackage{amsfonts}
\usepackage{amsmath}
\usepackage{amssymb}
\usepackage{graphicx}%
\setcounter{MaxMatrixCols}{30}
%TCIDATA{OutputFilter=latex2.dll}
%TCIDATA{Version=4.00.0.2321}
%TCIDATA{CSTFile=revtex4.cst}
%TCIDATA{Created=Wednesday, June 05, 2002 11:05:27}
%TCIDATA{LastRevised=Thursday, June 06, 2002 21:19:29}
%TCIDATA{<META NAME="GraphicsSave" CONTENT="32">}
%TCIDATA{<META NAME="DocumentShell" CONTENT="Articles\SW\REVTeX 4">}
%TCIDATA{Language=American English}
\newtheorem{theorem}{Theorem}
\newtheorem{acknowledgement}[theorem]{Acknowledgement}
\newtheorem{algorithm}[theorem]{Algorithm}
\newtheorem{axiom}[theorem]{Axiom}
\newtheorem{claim}[theorem]{Claim}
\newtheorem{conclusion}[theorem]{Conclusion}
\newtheorem{condition}[theorem]{Condition}
\newtheorem{conjecture}[theorem]{Conjecture}
\newtheorem{corollary}[theorem]{Corollary}
\newtheorem{criterion}[theorem]{Criterion}
\newtheorem{definition}[theorem]{Definition}
\newtheorem{example}[theorem]{Example}
\newtheorem{exercise}[theorem]{Exercise}
\newtheorem{lemma}[theorem]{Lemma}
\newtheorem{notation}[theorem]{Notation}
\newtheorem{problem}[theorem]{Problem}
\newtheorem{proposition}[theorem]{Proposition}
\newtheorem{remark}[theorem]{Remark}
\newtheorem{solution}[theorem]{Solution}
\newtheorem{summary}[theorem]{Summary}
\newenvironment{proof}[1][Proof]{\noindent\textbf{#1.} }{\ \rule{0.5em}{0.5em}}
\begin{document}
\preprint{ }
\title[Functionals]{Functionals}
\author{Brian Weiner}
\affiliation{PSU}
\date{06/06/2002}
\startpage{1}
\maketitle
\tableofcontents


\section{Functionals}

Let us compare functions, $f,l$, of discrete vectors to those, $F,L$, of
functions i.e. functions of vectors expressed in the basis $\left\{
\left\vert e_{j}\right\rangle \equiv\left\vert j\right\rangle ,\,j\in
\mathbb{N}\right\}  $ to those expressed in the basis $\left\{  \left\vert
t\right\rangle ,\,t\in\left[  0,\infty\right]  \right\}  $. The discrete case
is
\begin{align}
f\left(  \mathbf{x},n\right)   &  =\sum_{1\leq j\leq n}l\left(  x_{j}%
,j\right)  =\sum_{1\leq j\leq n}f_{0j}+\sum_{1\leq j\leq n}f_{1j}x_{j}%
+\sum_{1\leq j,k\leq n}f_{2jk}x_{j}x_{k}+\cdots\nonumber\\
&  =f_{0}+\mathbf{f}_{1}\cdot\mathbf{x+f}_{2}\cdot\mathbf{x\otimes x}%
+\cdots\nonumber\\
&  =f_{0}+\left\langle \mathbf{f}_{1}|\mathbf{x}\right\rangle +\left\langle
\mathbf{f}_{2}|\mathbf{x\otimes x}\right\rangle +\cdots
\end{align}
where%
\begin{equation}
l\left(  x_{j},j\right)  =f_{0j}+f_{1j}x_{j}+\sum_{1\leq k\leq n}f_{2jk}%
x_{j}x_{k}+\cdots
\end{equation}
We can take derivatives wrt the coordinates of the basis vectors $\left\{
\left\vert e_{j}\right\rangle \right\}  $ so that we have%
\begin{equation}
df\left(  \mathbf{x},n\right)  \left(  \delta\mathbf{x}\right)  =\sum_{1\leq
j\leq n}\frac{\partial f\left(  \mathbf{x},n\right)  }{\partial x_{j}}\delta
x_{j}%
\end{equation}
If the differential $df\left(  \mathbf{x},n\right)  $ exists then the
derivative is of the Fr\u{e}chet type
\begin{align}
\frac{\partial f\left(  \mathbf{x},n\right)  }{\partial x_{r}}  &
=f_{1r}+\sum_{1\leq j\leq n}\left\{  f_{2jr}+f_{2rj}\right\}  x_{j}%
+\cdots\nonumber\\
&  =\left\langle \mathbf{f}_{1}|e_{r}\right\rangle +\left\langle
\mathbf{f}_{2}|\mathbf{x\otimes}e_{r}+e_{r}\mathbf{\otimes x}\right\rangle
+\cdots\nonumber\\
\frac{\partial l\left(  x_{j},j\right)  }{\partial x_{r}}  &  =f_{1j}%
\delta_{jr}+f_{2rj}x_{j}+\sum_{1\leq k\leq n}f_{2jk}x_{k}\delta_{jr}+\cdots
\end{align}
and the consistency check gives
\begin{align}
\frac{\partial f\left(  \mathbf{x},n\right)  }{\partial x_{r}}  &
=\frac{\partial}{\partial x_{r}}\sum_{1\leq j\leq n}l\left(  x_{j},j\right)
=\sum_{1\leq j\leq n}\frac{\partial l\left(  x_{j},j\right)  }{\partial x_{r}%
}\nonumber\\
&  =f_{1r}+\sum_{1\leq j\leq n}\left\{  f_{2rj}+f_{2rj}\right\}  x_{j}+\cdots
\end{align}


Compare this to the continuous case
\begin{align}
&  F\left(  x,\tau\right)  =\int_{0}^{\tau}L\left(  x,t\right)  dt\nonumber\\
&  =\int_{0}^{\tau}F_{0}\left(  t\right)  dt+\int_{0}^{\tau}F_{1}\left(
t_{1}\right)  x\left(  t_{1}\right)  dt_{1}+\int_{0}^{\tau}F_{2}\left(
t_{1},t_{2}\right)  x\left(  t_{1}\right)  x\left(  t_{2}\right)  dt_{1}%
dt_{2}+\cdots\nonumber\\
&  =\int_{0}^{\tau}F_{0}\left(  t\right)  dt+\left\langle \mathbf{F}%
_{1}|\mathbf{x}\right\rangle +\left\langle \mathbf{F}_{2}|\mathbf{x\otimes
x}\right\rangle +\cdots\nonumber\\
&  L\left(  x,t\right)  =F_{0}\left(  t\right)  +F_{1}\left(  t\right)
x\left(  t\right)  +\int_{0}^{\tau}F_{2}\left(  t,t_{2}\right)  x\left(
t\right)  x\left(  t_{2}\right)  dt_{2}+\cdots
\end{align}
From the fundamental theorem of calculus
\begin{equation}
L\left(  x,\tau\right)  =\frac{\partial F\left(  x,\tau\right)  }{\partial
\tau}\Leftrightarrow L\left(  x,t\right)  =\frac{\partial F\left(  x,t\right)
}{\partial t}\forall t\in\left[  0,\tau\right]
\end{equation}
We can take derivatives wrt the coordinates of the basis vectors $\left\{
\left\vert t\right\rangle \right\}  $
\begin{align}
&  \frac{\delta F\left(  x,\tau\right)  }{\delta x\left(  t\right)  }%
=F_{1}\left(  t\right)  +\int_{0}^{\tau}\left\{  F_{2}\left(  t_{1},t\right)
+F_{2}\left(  t,t_{1}\right)  \right\}  x\left(  t_{1}\right)  dt_{1}%
+\cdots\nonumber\\
&  =\left\langle \mathbf{F}_{1}|t\right\rangle +\left\langle \mathbf{F}%
_{2}|\mathbf{x\otimes}t\right\rangle +\left\langle \mathbf{F}_{2}%
|t\mathbf{\otimes x}\right\rangle +\cdots\nonumber\\
&  dF\left(  x,\tau\right)  \left(  \delta x,\delta\tau\right)  =\int
_{0}^{\tau}\frac{\delta F\left(  x,\tau\right)  }{\delta x\left(  t\right)
}\delta x\left(  t\right)  dt+\frac{\partial F\left(  x,\tau\right)
}{\partial\tau}\delta\tau\nonumber\\
&  \frac{\delta L\left(  x,t^{\prime}\right)  }{\delta x\left(  t\right)
}=F_{1}\left(  t\right)  \delta\left(  t-t^{\prime}\right)  +F_{2}\left(
t,t^{\prime}\right)  x\left(  t^{\prime}\right)  +\int_{0}^{\tau}F_{2}\left(
t,t_{2}\right)  \delta\left(  t-t^{\prime}\right)  x\left(  t_{2}\right)
dt_{2}+\cdots
\end{align}
The time duration derivatives are
\begin{equation}
\frac{\partial F\left(  x,\tau\right)  }{\partial\tau}=\int_{0}^{\tau}%
\frac{\partial L\left(  x,t\right)  }{\partial t}dt=L\left(  x,\tau\right)
\end{equation}
when $L\left(  x,0\right)  =0.$ Which in detail gives
\begin{equation}
\frac{\partial L\left(  x,t\right)  }{\partial t}=\frac{\partial}{\partial
t}\left\{  F_{0}\left(  t\right)  +F_{1}\left(  t\right)  x\left(  t\right)
+\int_{0}^{\tau}F_{2}\left(  t,t_{2}\right)  x\left(  t\right)  x\left(
t_{2}\right)  dt_{2}+\cdots\right\}
\end{equation}
Clearly
\begin{equation}
\frac{\delta F\left(  x,\tau\right)  }{\delta x\left(  t\right)  }=\int
_{0}^{\tau}\frac{\delta L\left(  x,t^{\prime}\right)  }{\delta x\left(
t\right)  }dt^{\prime}%
\end{equation}
and
\begin{equation}
\frac{\partial}{\partial t}\left\{  \frac{\delta F\left(  x,\tau\right)
}{\delta x\left(  t\right)  }\right\}  =\frac{\delta L\left(  x,t^{\prime
}\right)  }{\delta x\left(  t\right)  }%
\end{equation}
The differential $dF$ of $F$ is then given by
\begin{align}
&  dF\left(  x,\tau\right)  \left(  \delta x,\delta\tau\right)  =\int
_{0}^{\tau}dL\left(  x,t\right)  \left(  \delta x,\delta t\right)
dt\nonumber\\
&  =\int_{0}^{\tau}\left\{  \frac{\partial L\left(  x,t\right)  }{\partial
x\left(  t\right)  }\delta x\left(  t\right)  +\frac{\partial L\left(
x,t\right)  }{\partial t}\delta t\right\}  dt\nonumber\\
&  =\int_{0}^{\tau}\frac{\partial L\left(  x,t\right)  }{\partial x\left(
t\right)  }\delta x\left(  t\right)  dt+\frac{\partial F\left(  x,\tau\right)
}{\partial\tau}\delta\tau\nonumber\\
&  =\int_{0}^{\tau}\frac{\partial L\left(  x,t\right)  }{\partial x\left(
t\right)  }\delta x\left(  t\right)  dt+L\left(  x,\tau\right)  \delta\tau
\end{align}
and
\begin{equation}
\delta^{\left(  1\right)  }F\left(  x,\tau\right)  =dF\left(  x,\tau\right)
\left(  \delta x,\delta\tau\right)
\end{equation}



\end{document}